\documentclass[11pt]{book}

\usepackage[utf8]{inputenc}
\usepackage[T1]{fontenc}
\usepackage[margin=1in]{geometry}
\usepackage{amsmath,amssymb,amsthm,mathtools}
\usepackage{booktabs}
\usepackage{hyperref}

\hypersetup{colorlinks=true,linkcolor=blue,citecolor=blue,urlcolor=blue}

\newcommand{\rd}{\mathrm{d}}
\newcommand{\pdv}[2]{\dfrac{\partial #1}{\partial #2}}
\newcommand{\fdv}[2]{\dfrac{\delta #1}{\delta #2}}
\newcommand{\Lag}{\mathcal{L}}
\newcommand{\Ham}{\mathcal{H}}

\theoremstyle{definition}
\newtheorem{definition}{Definition}[section]
\newtheorem{example}{Example}[section]
\theoremstyle{remark}
\newtheorem{remark}{Remark}[section]

\title{Two Roads to Infinite-Dimensional Dynamics}
\author{}
\date{}

\begin{document}

\chapter{Two Roads to Infinite-Dimensional Dynamics: Field Gradients as Momenta versus Many-Fingered Time}
\label{ch:infinite-dof}

\section{Introduction}

Ordinary Hamiltonian mechanics is built for systems with finitely many degrees of freedom: a configuration is a point $q=(q^1,\dots,q^N)$, a state is a point $(q,p)$ in a $2N$-dimensional phase space, and time is a single external parameter with respect to which everything evolves. The moment we pass to a \emph{field} --- an elastic string, the electromagnetic field, a wavefunctional --- the configuration is no longer a finite list of numbers but a whole function $\phi(\mathbf{x})$ on space, and the naive counting ``one degree of freedom per value of $\mathbf{x}$'' makes the number of degrees of freedom literally infinite (indeed continuously infinite).

Two logically independent strategies exist for extending the finite-dimensional Hamiltonian apparatus to this situation, and this chapter develops both in full technical detail.

\begin{itemize}
\item[(I)] \textbf{Field gradients as momenta.} Keep a single, distinguished time, but do not stop at promoting only the time-derivative $\dot\phi$ to a momentum: promote \emph{every} derivative of the field, including its spatial gradient $\nabla\phi$, to an independent momentum-like variable. This is the road to the \emph{De Donder--Weyl covariant Hamiltonian formalism}, in which the phase space attached to each spacetime point is finite-dimensional, and the spatial gradient literally \emph{becomes} (up to raising an index) a momentum component.

\item[(II)] \textbf{Coordinates and multidimensional time.} Keep the ordinary field and its ordinary canonical momentum density, but do not stop at a single external time: give \emph{every point of space its own clock}. This is the road to Dirac's many-time formalism, its continuum avatar the \emph{Tomonaga--Schwinger equation}, and its fully geometric incarnation in the \emph{many-fingered time} of canonical general relativity.
\end{itemize}

The two roads generalize \emph{different} ingredients of finite-dimensional mechanics: (I) enlarges the momentum content while leaving time alone; (II) enlarges the time content while leaving the momentum alone. Section~\ref{sec:synthesis} makes this complementarity precise. Both are developed here from the same starting point --- the continuum limit of a chain of coupled oscillators --- so that their relation to one another, and to the ordinary canonical field theory that sits between them, is transparent.

\section{Notation and Conventions}

Latin indices $i,j,k,\dots$ run over spatial values ($1,2,3$, or just $1$ for the one-dimensional string of \S\ref{sec:continuum-limit}); Greek indices $\mu,\nu,\dots$ run over spacetime values $0,1,2,3$. Repeated indices are summed (Einstein convention). The spacetime metric is $\eta_{\mu\nu}=\mathrm{diag}(+1,-1,-1,-1)$, so $\partial^\mu=\eta^{\mu\nu}\partial_\nu$ gives $\partial^0=\partial_0=\partial/\partial t$ and $\partial^i=-\partial_i$. We set $\hbar=c=1$ in the relativistic sections and restore them where it clarifies a non-relativistic formula. $\nabla\phi$ denotes the ordinary spatial gradient $(\partial_1\phi,\partial_2\phi,\partial_3\phi)$.

\begin{definition}[Functional derivative]
\label{def:funcderiv}
For a functional $F[\phi]=\int \rd^3x\, f\big(\phi(\mathbf{x}),\nabla\phi(\mathbf{x})\big)$ of a field $\phi$ on space, the functional derivative $\delta F/\delta\phi(\mathbf{x})$ is the distribution defined by
\begin{equation}
\delta F = \int \rd^3x\; \fdv{F}{\phi(\mathbf{x})}\,\delta\phi(\mathbf{x}) + O(\delta\phi^2),
\end{equation}
for arbitrary variations $\delta\phi$ vanishing at spatial infinity. Concretely,
\begin{equation}
\fdv{F}{\phi(\mathbf{x})} = \pdv{f}{\phi} - \nabla\cdot\pdv{f}{(\nabla\phi)}.
\end{equation}
\end{definition}

\section{From a Finite Chain to a Continuous Field}
\label{sec:continuum-limit}

Both approaches below start from the same elementary construction, so it is worth doing once, carefully. Consider $N$ point masses $m$ threaded on a taut string of tension $\tau$, spaced a distance $a$ apart, with transverse displacements $\eta_i(t)$. For small displacements the Lagrangian is
\begin{equation}
L = \sum_{i} \left[\frac{m}{2}\dot\eta_i^2 - \frac{\tau}{2a}(\eta_{i+1}-\eta_i)^2\right].
\end{equation}
Each $\eta_i$ is an independent generalized coordinate; the system has $N$ degrees of freedom, and ordinary Hamiltonian mechanics applies with $p_i=\partial L/\partial\dot\eta_i = m\dot\eta_i$.

Now take the continuum limit: $a\to0$, $N\to\infty$, with the linear mass density $\mu \equiv m/a$ and the tension $\tau$ held fixed, and relabel $\eta_i(t)\to \phi(x,t)$ with $x=ia$. Using $(\eta_{i+1}-\eta_i)/a \to \partial\phi/\partial x$ and $\sum_i a \to \int \rd x$,
\begin{equation}
L \;\longrightarrow\; \int \rd x\, \Lag(\phi,\dot\phi,\phi'), \qquad
\Lag = \frac{\mu}{2}\dot\phi^2 - \frac{\tau}{2}\phi'^2 ,
\label{eq:string-lagrangian}
\end{equation}
where $\phi'\equiv\partial\phi/\partial x$. Two things have happened simultaneously. First, the discrete index $i$ has become a continuous label $x$: the ``coordinate'' is now the whole function $\phi(\cdot,t)$, an element of an infinite-dimensional configuration space. Second, the nearest-neighbor coupling term, which in the discrete system was an ordinary potential energy coupling adjacent coordinates, has turned into a term built from the \emph{spatial gradient} $\phi'$ of the field.

This second fact is the seed of Approach~I: the gradient $\phi'$ is, dimensionally and structurally, on exactly the same footing as the time derivative $\dot\phi$ --- both are derivatives of $\phi$ with respect to one of the independent variables $(t,x)$ that label a spacetime point. The only reason the Lagrangian formalism treats them differently is that the least-action principle singles out $t$ as ``the'' evolution parameter. Sections~\ref{sec:canonical} and \ref{sec:dw} explore what happens when that asymmetry is removed.

The same discretization, read differently, is the seed of Approach~II: before taking $N\to\infty$, one could equally well have given \emph{each} $\eta_i$ its own independent time parameter $t_i$, since nothing in Hamilton's principle for a many-particle system requires a single common time. Section~\ref{sec:dirac-multitime} develops this idea for particles, and \S\ref{sec:tomonaga-schwinger}--\S\ref{sec:dirac-algebra} take the same $N\to\infty$ limit \emph{on the time side} rather than the coordinate side.

\section{Approach I: Field Gradients as Momenta}
\label{sec:approach1}

\subsection{The Instant-Form Canonical Formalism}
\label{sec:canonical}

The standard construction singles out $t$: the canonical momentum \emph{density} conjugate to $\phi(\mathbf{x})$ is
\begin{equation}
\pi(\mathbf{x},t) \equiv \fdv{L}{\dot\phi(\mathbf{x})} = \pdv{\Lag}{\dot\phi}\bigg|_{\mathbf{x}}.
\end{equation}
For the string \eqref{eq:string-lagrangian}, $\pi = \mu\dot\phi$. The Hamiltonian density follows from the ordinary (single-variable) Legendre transform in $\dot\phi$ alone:
\begin{equation}
\Ham = \pi\dot\phi - \Lag = \frac{\pi^2}{2\mu} + \frac{\tau}{2}\phi'^2 ,
\qquad
H[\phi,\pi] = \int \rd x\, \Ham .
\end{equation}
Phase space is now the (infinite-dimensional) space of pairs of functions $(\phi(\cdot),\pi(\cdot))$, equipped with the canonical Poisson brackets
\begin{equation}
\{\phi(x),\pi(y)\} = \delta(x-y), \qquad \{\phi(x),\phi(y)\}=\{\pi(x),\pi(y)\}=0.
\label{eq:canonical-pb}
\end{equation}
Hamilton's equations become functional-derivative equations,
\begin{equation}
\dot\phi(x) = \fdv{H}{\pi(x)} = \frac{\pi(x)}{\mu}, \qquad
\dot\pi(x) = -\fdv{H}{\phi(x)} = \tau\,\phi''(x),
\end{equation}
using Definition~\ref{def:funcderiv}, and combine into the wave equation $\mu\ddot\phi=\tau\phi''$, wave speed $c^2=\tau/\mu$, as expected.

Notice the asymmetry that survives this construction: $\dot\phi$ was promoted to a momentum $\pi$, but $\phi'$ was not --- it simply sits inside $\Ham$ as a function of the coordinate field. Space and time, symmetric in the original wave equation, are treated unequally by the canonical formalism, precisely because the formalism is built around a choice of ``the'' time direction.

\subsection{Restoring Space--Time Symmetry: The De Donder--Weyl Formalism}
\label{sec:dw}

The De Donder--Weyl (DW) formalism removes this asymmetry by performing a Legendre transform with respect to \emph{all} derivatives of the field at once, rather than the time derivative alone. Let $\Lag=\Lag(\phi,\partial_\mu\phi,x^\mu)$ be a Lagrangian \emph{density on spacetime} (not merely on space). Define, for every spacetime direction $\mu$, a \emph{polymomentum}
\begin{equation}
\pi^\mu \equiv \pdv{\Lag}{(\partial_\mu\phi)}.
\end{equation}
There is one such momentum for the time derivative ($\mu=0$, numerically the same as the canonical $\pi$ above) \emph{and one for each spatial derivative} ($\mu=i$): the spatial polymomenta $\pi^i$ are conjugate to the components of the gradient $\partial_i\phi$. The covariant (De Donder--Weyl) Hamiltonian function is defined by the corresponding multivariable Legendre transform,
\begin{equation}
\Ham_{\mathrm{DW}}(\phi,\pi^\mu,x^\mu) \equiv \pi^\mu\,\partial_\mu\phi - \Lag ,
\label{eq:dw-hamiltonian}
\end{equation}
in which $\partial_\mu\phi$ on the right is understood as the function of $(\phi,\pi^\mu)$ obtained by inverting the definitions of the $\pi^\mu$ (this inversion is possible whenever the Lagrangian is \emph{regular}, i.e.\ the Hessian $\partial^2\Lag/\partial(\partial_\mu\phi)\partial(\partial_\nu\phi)$ is nondegenerate; see \S\ref{sec:scope} for the gauge-theory exception). The De Donder--Weyl field equations, equivalent to the Euler--Lagrange equation for $\Lag$, are
\begin{equation}
\partial_\mu\phi = \pdv{\Ham_{\mathrm{DW}}}{\pi^\mu} \quad(\text{no sum}), \qquad
\partial_\mu\pi^\mu = -\pdv{\Ham_{\mathrm{DW}}}{\phi} \quad(\text{sum over }\mu).
\label{eq:dw-eom}
\end{equation}
The first set simply reproduces, for every $\mu$, the definition of $\pi^\mu$ (exactly as $\dot q=\partial H/\partial p$ reproduces $p=m\dot q$ in point mechanics); the second is a genuine equation of motion, a spacetime divergence condition on the polymomentum vector $\pi^\mu$.

The essential structural point is this: at a fixed spacetime point, the pair $(\phi,\pi^\mu)$ lives in a space of dimension $1+(d+1)$ for a spacetime of $d$ spatial dimensions --- \emph{finite}, unlike the canonical phase space of \S\ref{sec:canonical}, which is a space of functions on an entire spatial slice. The price paid for this finite-dimensionality is that the DW ``Hamiltonian'' no longer generates evolution in a single external time; it is a function \emph{on spacetime}, and \eqref{eq:dw-eom} is a first-order system of \emph{partial}, not ordinary, differential equations. This is the sense in which DW theory is a genuinely \emph{covariant} Hamiltonian formalism: it has traded a single infinite-dimensional phase space evolving in time for a finite-dimensional ``polyphase space'' fibered over all of spacetime.

\subsection{Worked Example: The Free Klein--Gordon Field}
\label{sec:kg-example}

Let $\phi$ be a real scalar with
\begin{equation}
\Lag = \frac{1}{2}\eta^{\mu\nu}\partial_\mu\phi\,\partial_\nu\phi - \frac{1}{2}m^2\phi^2
     = \frac{1}{2}\dot\phi^2 - \frac{1}{2}(\nabla\phi)^2 - \frac{1}{2}m^2\phi^2 .
\end{equation}

\paragraph{Canonical formalism.} $\pi=\dot\phi$, and
\begin{equation}
\Ham = \frac{1}{2}\pi^2 + \frac{1}{2}(\nabla\phi)^2 + \frac{1}{2}m^2\phi^2 ,
\end{equation}
manifestly non-negative: this is the ordinary energy density, the $T^{00}$ component of the stress tensor, and it is frame-dependent (it refers to a specific choice of time direction).

\paragraph{De Donder--Weyl formalism.} The polymomenta are $\pi^\mu = \partial^\mu\phi$, i.e.
\begin{equation}
\pi^0 = \partial^0\phi = \dot\phi, \qquad \pi^i = \partial^i\phi = -\partial_i\phi .
\label{eq:kg-polymomenta}
\end{equation}
Equation \eqref{eq:kg-polymomenta} is the precise, worked-example content of ``field gradients as momenta'': the spatial gradient $\partial_i\phi$ \emph{is}, up to the sign forced by the metric, the spatial polymomentum $\pi^i$. Carrying out the Legendre transform \eqref{eq:dw-hamiltonian},
\begin{align}
\Ham_{\mathrm{DW}} &= \pi^\mu\partial_\mu\phi - \Lag
= \big[(\pi^0)^2 - |\boldsymbol\pi|^2\big] - \left[\frac12(\pi^0)^2-\frac12|\boldsymbol\pi|^2-\frac12 m^2\phi^2\right] \notag\\[2pt]
&= \frac{1}{2}(\pi^0)^2 - \frac{1}{2}|\boldsymbol\pi|^2 + \frac{1}{2}m^2\phi^2
= \frac{1}{2}\pi^\mu\pi_\mu + \frac{1}{2}m^2\phi^2 ,
\label{eq:kg-dw-hamiltonian}
\end{align}
where $\boldsymbol\pi=(\pi^1,\pi^2,\pi^3)$. Unlike the canonical $\Ham$, this is a \emph{Lorentz scalar} (it is built from $\pi^\mu\pi_\mu$, a full contraction), and it is \emph{not} sign-definite --- it is not an energy density at all, but a genuinely different, frame-independent object. The DW field equations \eqref{eq:dw-eom} give, trivially, $\partial_0\phi=\pi^0$ and $\partial_i\phi=-\pi^i$ (recovering \eqref{eq:kg-polymomenta}), and, non-trivially,
\begin{equation}
\partial_\mu\pi^\mu = -m^2\phi
\;\;\Longrightarrow\;\;
\ddot\phi - \nabla^2\phi = -m^2\phi
\;\;\Longrightarrow\;\;
\ddot\phi-\nabla^2\phi+m^2\phi = 0,
\end{equation}
the Klein--Gordon equation, confirming that \eqref{eq:kg-dw-hamiltonian} generates the correct dynamics.

\begin{table}[h]
\centering
\begin{tabular}{@{}lll@{}}
\toprule
 & Canonical (instant form) & De Donder--Weyl \\
\midrule
Momentum conjugate to $\phi$ & $\pi=\dot\phi$ only & $\pi^\mu=\partial^\mu\phi$, all $\mu$ \\
Role of $\nabla\phi$ & appears in $\Ham$, not a momentum & $\nabla\phi = -\boldsymbol\pi$, a momentum \\
Phase space (per instant / per point) & $\infty$-dimensional (functions on $\Sigma$) & finite-dimensional \\
Character of $\Ham$ & energy density, frame-dependent & Lorentz scalar, sign-indefinite \\
Evolution equations & functional ODEs in $t$ & PDEs on spacetime \\
\bottomrule
\end{tabular}
\caption{The Klein--Gordon field in the two Hamiltonian formalisms of Approach I.}
\end{table}

\subsection{The Multisymplectic Phase Space}

The geometric home of the DW formalism is the \emph{multisymplectic manifold}: instead of the symplectic $2$-form $\rd q\wedge \rd p$ of point mechanics, one equips the finite-dimensional space of variables $(x^\mu,\phi,\pi^\mu)$ with a $(d+2)$-form
\begin{equation}
\Omega = \rd\phi\wedge \rd\pi^\mu \wedge \rd^{d}x_\mu ,
\end{equation}
where $\rd^d x_\mu$ denotes the volume form on spacetime with the $\mu$-th direction omitted. Just as Hamilton's equations in point mechanics are the statement that the flow generated by $H$ preserves the symplectic form, the DW field equations \eqref{eq:dw-eom} are the statement that the ``polysymplectic'' structure $\Omega$ is preserved by the field configuration $(\phi(x),\pi^\mu(x))$. We will not develop the jet-bundle machinery needed to make this fully precise (see the Further Reading), but the qualitative moral is the one already drawn from the worked example: DW theory replaces an infinite-dimensional phase space evolving in one time with a finite-dimensional fiber bundle sitting over all of spacetime, at the cost of losing a single privileged notion of ``simultaneous state.''

\subsection{Scope and Limitations}
\label{sec:scope}

The covariant Legendre transform \eqref{eq:dw-hamiltonian} requires the map $\partial_\mu\phi \mapsto \pi^\mu$ to be invertible. This fails for gauge theories: for the free electromagnetic field, $\Lag=-\tfrac14F_{\mu\nu}F^{\mu\nu}$, the polymomentum $\pi^{\mu\nu}=\partial\Lag/\partial(\partial_\mu A_\nu)=F^{\mu\nu}$ conjugate to $\partial_0A_0$ vanishes identically, reflecting the gauge freedom $A_\mu\to A_\mu+\partial_\mu\lambda$ (Exercise~1 below works through this). Constrained/gauge DW theories exist but require an analogue of the Dirac--Bergmann procedure for constrained systems, generalized to the covariant setting; we do not pursue this here.

\section{Approach II: Coordinates and Multidimensional Time}
\label{sec:approach2}

\subsection{Dirac's Many-Time Formalism for Particles}
\label{sec:dirac-multitime}

Consider $N$ non-relativistic particles with the ordinary (single-time) Schr\"odinger equation $i\,\partial_t\psi = H\psi$, $H=\sum_a H_a$. Nothing in principle forbids giving each particle its own time coordinate: define $\psi(x_1,t_1;\dots;x_N,t_N)$, evolving according to
\begin{equation}
i\,\pdv{\psi}{t_a} = H_a\,\psi , \qquad a=1,\dots,N,
\label{eq:multitime-schrodinger}
\end{equation}
with each $H_a$ a function of the coordinates/momenta of particle $a$ alone (and possibly of the other particles' \emph{coordinates}, but not their momenta or times). The physical wavefunction is recovered on the ``diagonal'' $t_1=\dots=t_N\equiv t$.

For \eqref{eq:multitime-schrodinger} to define $\psi$ consistently --- so that the value obtained at a point $(t_1,\dots,t_N)$ does not depend on the order in which the $N$ equations are integrated to get there --- the mixed partial derivatives must agree:
\begin{equation}
\pdv{}{t_b}\left(i\pdv{\psi}{t_a}\right) = \pdv{}{t_a}\left(i\pdv{\psi}{t_b}\right)
\;\;\Longrightarrow\;\;
H_aH_b\psi = H_bH_a\psi \quad \text{for all }\psi,
\end{equation}
i.e.\ the \textbf{Dirac integrability condition}
\begin{equation}
[H_a,H_b]=0, \qquad a\neq b.
\label{eq:dirac-integrability}
\end{equation}
This is the precise sense in which the times $t_a$ are ``independent coordinates'': the system \eqref{eq:multitime-schrodinger} defines a genuine, path-independent evolution on the $N$-dimensional space of times $(t_1,\dots,t_N)$ if and only if \eqref{eq:dirac-integrability} holds. For free or suitably decoupled particles this holds trivially ($[H_a,H_b]=0$ when $H_a,H_b$ act on different tensor factors); for particles interacting through an instantaneous potential $V(x_a-x_b)$ it generally \emph{fails} (Exercise~3), which is one of the classic obstructions to a naive multi-time formulation of interacting particle mechanics, relativistic or not.

\subsection{The Continuum Limit: Tomonaga--Schwinger Theory}
\label{sec:tomonaga-schwinger}

Now take the same $N\to\infty$ limit that produced a field in \S\ref{sec:continuum-limit}, but apply it to the \emph{time} index of \S\ref{sec:dirac-multitime} rather than to a coordinate: replace the discrete particle label $a$ by a continuous label $\mathbf{x}$ ranging over a spacelike Cauchy hypersurface $\Sigma$, and replace the $N$ independent times $t_a$ by a \emph{function} $\tau(\mathbf{x})$ describing, at every point of $\Sigma$, how far that point has been displaced along the local timelike normal. A ``moment of time'' is no longer a single number $t$; it is an entire spacelike hypersurface $\Sigma$, and the surface can be advanced by different amounts at different points --- this is the continuum incarnation of ``multidimensional time.''

The state is a wave \emph{functional} $\Psi[\phi(\cdot);\Sigma]$ of the field configuration on $\Sigma$ and of $\Sigma$ itself. Under an infinitesimal, purely local deformation of $\Sigma$ by a normal displacement $\delta\tau(\mathbf{x})$, the \textbf{Tomonaga--Schwinger equation} states
\begin{equation}
i\,\fdv{\Psi[\Sigma]}{\tau(\mathbf{x})} = \Ham(\mathbf{x})\,\Psi[\Sigma],
\label{eq:tomonaga-schwinger}
\end{equation}
where $\Ham(\mathbf{x})$ is the Hamiltonian density (operator) at the spacetime point corresponding to $\mathbf{x}\in\Sigma$. Equation \eqref{eq:tomonaga-schwinger} is the direct continuum analogue of \eqref{eq:multitime-schrodinger}: it says that the surface can be bent independently at every point, with the local Hamiltonian density $\Ham(\mathbf{x})$ playing exactly the role that $H_a$ played for particle $a$.

Exactly as before, consistency (independence of $\Psi[\Sigma_{\mathrm{final}}]$ from the sequence of local deformations used to reach $\Sigma_{\mathrm{final}}$ from some reference $\Sigma_0$) requires the continuum integrability condition
\begin{equation}
[\Ham(\mathbf{x}),\Ham(\mathbf{y})] = 0 \qquad \text{for all } \mathbf{x}\neq\mathbf{y}\text{ on }\Sigma.
\label{eq:microcausality}
\end{equation}
Because $\Sigma$ is spacelike, any two of its points are spacelike separated; \eqref{eq:microcausality} is exactly the statement of \textbf{microcausality} in local quantum field theory --- that operators built from the fields at spacelike-separated points commute. For the free scalar field this is a theorem, not an assumption: the field commutator is the Pauli--Jordan function,
\begin{equation}
[\phi(x),\phi(y)] = i\,\Delta(x-y), \qquad \Delta(x-y)=0 \ \text{for } (x-y)^2<0,
\end{equation}
and since $\Ham(\mathbf{x})$ is built locally from $\phi$ and its derivatives at $\mathbf{x}$, \eqref{eq:microcausality} follows. Microcausality is thus the microscopic reason a consistent many-fingered-time description of quantum field dynamics exists at all; it is also, not coincidentally, the same ingredient responsible for the Lorentz invariance of the resulting S-matrix in interacting theories.

\subsection{Geometrizing Many-Fingered Time: Embeddings and Deformations}
\label{sec:embeddings}

To make ``the space of all possible clocks at all points of space'' into a genuine dynamical arena, following Dirac (1951) it is useful to promote the hypersurface $\Sigma$ itself to a dynamical object. Fix a background spacetime with global coordinates $X^\mu$ and metric $\eta_{\mu\nu}$, and describe $\Sigma$ by an \textbf{embedding} $X^\mu=X^\mu(\mathbf{x})$, where $\mathbf{x}=(x^1,x^2,x^3)$ are coordinates on an abstract label space (``space'' in the intrinsic sense). The induced metric on $\Sigma$ is
\begin{equation}
h_{ij}(\mathbf{x}) = \partial_iX^\mu\,\partial_jX^\nu\,\eta_{\mu\nu},
\end{equation}
and there is a unique future-pointing unit normal $n^\mu(\mathbf{x})$ with $n\cdot n=1$, $n_\mu\partial_iX^\mu=0$. Treating $X^\mu(\mathbf{x})$ as canonical coordinates with conjugate momenta $P_\mu(\mathbf{x})$, and letting $\pi(\mathbf{x})$ be the canonical momentum of whatever matter field $\phi$ lives on $\Sigma$, the dynamics is entirely encoded in a pair of \textbf{constraints} at every point $\mathbf{x}$,
\begin{align}
\Ham_\perp(\mathbf{x}) &\equiv n^\mu(\mathbf{x})\,P_\mu(\mathbf{x}) - \Ham_{\mathrm{matter}}(\mathbf{x}) \;\approx\; 0, \label{eq:superham}\\
\Ham_i(\mathbf{x}) &\equiv \partial_iX^\mu(\mathbf{x})\,P_\mu(\mathbf{x}) + \pi(\mathbf{x})\,\partial_i\phi(\mathbf{x}) \;\approx\; 0, \label{eq:supermomentum}
\end{align}
where $\Ham_{\mathrm{matter}}(\mathbf{x})$ is the ordinary canonical Hamiltonian density of \S\ref{sec:canonical}, re-expressed using the induced metric $h_{ij}$. Here $\Ham_\perp$ (the \emph{super-Hamiltonian}) generates a purely normal deformation of $\Sigma$ at $\mathbf{x}$ --- an advance of the local clock --- while $\Ham_i$ (the \emph{super-momentum}) generates a purely tangential deformation, i.e.\ a relabeling of the points of $\Sigma$ that leaves the surface itself fixed. The total generator of an arbitrary infinitesimal deformation, specified by a \emph{lapse} $N(\mathbf{x})$ (normal displacement) and a \emph{shift} $N^i(\mathbf{x})$ (tangential displacement), is
\begin{equation}
H[N,N^i] = \int \rd^3x\, \big[N(\mathbf{x})\,\Ham_\perp(\mathbf{x}) + N^i(\mathbf{x})\,\Ham_i(\mathbf{x})\big].
\end{equation}
Because $N(\mathbf{x})$ and $N^i(\mathbf{x})$ are freely specifiable functions --- there is no preferred rate at which different points of space should advance in time --- the constraints $\Ham_\perp\approx0$, $\Ham_i\approx0$ must hold identically: this is exactly the statement, familiar from the ADM formulation of general relativity, that the lapse and shift are non-dynamical, and that ``time'' in field theory (and in gravity, once the embedding is allowed to depend on a genuinely dynamical metric rather than a fixed background) is not a single external parameter but a free function of $\mathbf{x}$ --- literally an infinite collection of local clock rates, one ``finger'' of time per spatial point.

\subsection{The Dirac--Teitelboim Hypersurface-Deformation Algebra}
\label{sec:dirac-algebra}

The consistency of this picture --- that deforming $\Sigma$ first by a lapse $N_1$ and then by $N_2$ gives the same final surface as the reverse order, up to a deformation generated by some \emph{other} combination of constraints --- is governed by the Poisson-bracket algebra of $\Ham_\perp$ and $\Ham_i$. Because the constraints \eqref{eq:superham}--\eqref{eq:supermomentum} are built purely from the kinematics of how a hypersurface can be embedded and deformed, their algebra is \emph{universal}: it holds for any matter content, and (Dirac 1951; Teitelboim 1973) takes the form
\begin{align}
\{\Ham_\perp(\mathbf{x}),\Ham_\perp(\mathbf{y})\} &= h^{ij}(\mathbf{x})\,\Ham_i(\mathbf{x})\,\partial_j\delta(\mathbf{x}-\mathbf{y}) - (\mathbf{x}\leftrightarrow\mathbf{y}), \label{eq:algebra1}\\
\{\Ham_\perp(\mathbf{x}),\Ham_i(\mathbf{y})\} &= \Ham_\perp(\mathbf{y})\,\partial_i\delta(\mathbf{x}-\mathbf{y}), \label{eq:algebra2}\\
\{\Ham_i(\mathbf{x}),\Ham_j(\mathbf{y})\} &= \Ham_i(\mathbf{y})\,\partial_j\delta(\mathbf{x}-\mathbf{y}) + \Ham_j(\mathbf{x})\,\partial_i\delta(\mathbf{x}-\mathbf{y}). \label{eq:algebra3}
\end{align}
Equation \eqref{eq:algebra3} alone is the ordinary Lie algebra of spatial diffeomorphisms (the Witt algebra in field-theory guise; Exercise~2 asks for its direct verification). Equations \eqref{eq:algebra1}--\eqref{eq:algebra2} are the genuinely new content: they say that two successive local time-advances do \emph{not} commute, but produce a spatial diffeomorphism, with a metric-dependent ``structure function'' $h^{ij}(\mathbf{x})$ rather than a structure \emph{constant}. Because the right-hand side of \eqref{eq:algebra1} depends on the dynamical field $h^{ij}$, this is not a genuine Lie algebra but what is sometimes called an \emph{open} or \emph{soft} algebra --- a fact of central importance in canonical quantum gravity, where it underlies the notorious ``problem of time'' in the Wheeler--DeWitt equation. We state \eqref{eq:algebra1}--\eqref{eq:algebra3} here without a full derivation (which requires working out how $n^\mu$ and $h_{ij}$ depend on $\partial_iX^\mu$ in detail); the key qualitative fact for this chapter is that the algebra is fixed by the pure geometry of embeddings, independent of which matter fields happen to live on $\Sigma$.

\subsection{Worked Example Revisited: Microcausality of the Free Scalar Field}

It is worth returning to \S\ref{sec:kg-example}'s Klein--Gordon field to see many-fingered time at its simplest: restrict to \emph{flat}, parallel slices, $X^0(\mathbf{x})=t$, $X^i(\mathbf{x})=x^i$, deformed only by a lapse $N(\mathbf{x})$ (no shift). Then $n^\mu=(1,0,0,0)$, $h_{ij}=\delta_{ij}$, and $\Ham_\perp(\mathbf{x})$ reduces exactly to the ordinary canonical energy density of \S\ref{sec:kg-example},
\begin{equation}
\Ham_\perp(\mathbf{x}) \to \Ham(\mathbf{x}) = \frac{1}{2}\pi(\mathbf{x})^2 + \frac{1}{2}(\nabla\phi(\mathbf{x}))^2 + \frac{1}{2}m^2\phi(\mathbf{x})^2 .
\end{equation}
A spatially constant lapse, $N(\mathbf{x})=N_0$, generates the ordinary global Hamiltonian $H=N_0\int\rd^3x\,\Ham(\mathbf{x})$, and every point of space advances together: single-time evolution is recovered as the special, ``rigid'' case of many-fingered time (Exercise~4). A spatially \emph{varying} lapse instead advances different regions of the field at different local rates while remaining on a family of parallel flat slices; \eqref{eq:microcausality}, i.e.\ $[\Ham(\mathbf{x}),\Ham(\mathbf{y})]=0$ for $\mathbf{x}\neq\mathbf{y}$, guarantees this is done consistently, and traces back, as shown in \S\ref{sec:tomonaga-schwinger}, to the vanishing of the Pauli--Jordan function outside the light cone.

\section{Synthesis: Two Orthogonal Generalizations}
\label{sec:synthesis}

\begin{table}[h]
\centering
\begin{tabular}{@{}lll@{}}
\toprule
 & Approach I (De Donder--Weyl) & Approach II (many-fingered time) \\
\midrule
What is generalized & momentum content & time content \\
Time & single, external & one local clock per spatial point \\
Momentum & one per spacetime direction $\mu$ & ordinary canonical $\pi(\mathbf{x})$ only \\
Phase space per point/instant & finite-dimensional & infinite-dimensional (functional) \\
Governing structure & multisymplectic geometry & hypersurface-deformation algebra \\
Consistency condition & regularity of Legendre transform & $[\Ham(\mathbf{x}),\Ham(\mathbf{y})]=0$, $\mathbf{x}\neq\mathbf{y}$ \\
\bottomrule
\end{tabular}
\caption{The two generalizations of finite-dimensional Hamiltonian mechanics compared.}
\end{table}

The two roads never actually meet in this chapter, and that is the point: Approach~I leaves the number of independent evolution parameters at one, and instead multiplies the momenta conjugate to a single field, so that spatial gradients acquire full momentum status; Approach~II leaves the momentum content exactly as in ordinary canonical field theory, and instead multiplies the number of independent evolution parameters, one per spatial point, so that ``the present moment'' becomes an entire freely-deformable hypersurface rather than a single number. Both are legitimate, and inequivalent, generalizations of the same finite-dimensional starting point $(q,p,t)$.

It is worth noting where a genuinely \emph{different} generalization of time sits relative to Approach~II. In Stueckelberg--Horwitz--Piron theory, dynamics is driven by a \emph{single} additional evolution parameter, an invariant ``historical time'' $\tau$, distinct from (and not to be confused with) the ordinary coordinate time $x^0$ that appears as one of the dynamical variables $x^\mu(\tau)$. This is neither Approach~I nor Approach~II in the sense developed here: it is not a multiplication of momenta (there is still one momentum $p^\mu$ per coordinate $x^\mu$), and it is not a multiplication of time into a function over space (there is exactly one $\tau$, shared by every degree of freedom, rather than one clock per spatial point). It occupies a third position in the taxonomy --- a single \emph{extra} time added on top of an unmodified spacetime, rather than either the spacetime derivatives being reorganized into momenta (Approach~I) or the single time being spread out over space into infinitely many local times (Approach~II).

\section*{Exercises}
\addcontentsline{toc}{section}{Exercises}

\begin{enumerate}
\item Derive the De Donder--Weyl Hamiltonian and covariant field equations for the free electromagnetic field, $\Lag=-\tfrac14F_{\mu\nu}F^{\mu\nu}$. Identify the polymomenta $\pi^{\mu\nu}=\partial\Lag/\partial(\partial_\mu A_\nu)$, show that $\pi^{00}=0$ identically, and relate this to the gauge freedom $A_\mu\to A_\mu+\partial_\mu\lambda$ and to the singularity of the covariant Legendre transform discussed in \S\ref{sec:scope}.

\item For the free scalar field, with $\Ham_i(\mathbf{x})=\pi(\mathbf{x})\partial_i\phi(\mathbf{x})$, compute $\{\Ham_i(\mathbf{x}),\Ham_j(\mathbf{y})\}$ directly from the canonical brackets \eqref{eq:canonical-pb} and check it against \eqref{eq:algebra3}.

\item Consider two interacting particles in Dirac's multi-time formalism, $H_1=p_1^2/2m_1+V(x_1-x_2)$, $H_2=p_2^2/2m_2+V(x_1-x_2)$. Show that $[H_1,H_2]\neq0$ in general, and discuss physically why an instantaneous potential is incompatible with a consistent multi-time description, and what kind of interaction (mediated, retarded) would restore \eqref{eq:dirac-integrability}.

\item Show explicitly that a family of parallel flat slices deformed by a spatially constant lapse $N(\mathbf{x})=N_0$ reduces \eqref{eq:superham}--\eqref{eq:supermomentum} to ordinary single-time Hamiltonian evolution with Hamiltonian $N_0 H$, as asserted at the end of \S\ref{sec:dirac-algebra}.
\end{enumerate}

\section*{Further Reading}
\addcontentsline{toc}{section}{Further Reading}

\begin{itemize}
\item H.~Goldstein, C.~Poole, J.~Safko, \textit{Classical Mechanics}, 3rd ed., ch.~13, for the discrete-chain-to-field continuum limit of \S\ref{sec:continuum-limit}.
\item P.A.M.~Dirac, ``Relativistic quantum mechanics with an application to Compton scattering'' and the multi-time formalism of \emph{Can.\ J.\ Math.} \textbf{3} (1951) 1, for \S\ref{sec:dirac-multitime} and \S\ref{sec:embeddings}.
\item S.~Tomonaga, \emph{Prog.\ Theor.\ Phys.} \textbf{1} (1946) 27; J.~Schwinger, \emph{Phys.\ Rev.} \textbf{74} (1948) 1439, for the equation of \S\ref{sec:tomonaga-schwinger}.
\item C.~Teitelboim, ``How commutators of constraints reflect the spacetime structure,'' \emph{Ann.\ Phys.} \textbf{79} (1973) 542, and K.~Kucha\v{r}, ``Geometrodynamics regained,'' for \S\ref{sec:dirac-algebra}.
\item C.J.~Isham, K.~Kucha\v{r}, ``Representations of spacetime diffeomorphisms. I, II,'' \emph{Ann.\ Phys.} \textbf{164} (1985), for the embedding-variable (parametrized field theory) construction of \S\ref{sec:embeddings}.
\item I.V.~Kanatchikov, and J.~Struckmeier \& H.~Reichau, on De Donder--Weyl / precanonical field theory, for \S\ref{sec:dw}--\S\ref{sec:scope}.
\end{itemize}

\end{document}
